%%
%% Copyright 2025 Oxford University Press
%%
%% This file is part of the 'oup-authoring-template Bundle'.
%% ---------------------------------------------
%%
%% It may be distributed under the conditions of the LaTeX Project Public
%% License, either version 1.3 of this license or (at your option) any
%% later version.  The latest version of this license is in
%%    https://www.latex-project.org/lppl.txt
%% and version 1.2 or later is part of all distributions of LaTeX
%% version 1999/12/01 or later.
%%
%% The list of all files belonging to the 'oup-authoring-template Bundle' is
%% given in the file `manifest.txt'.
%%
%% Template article for Oxford University Press's document class `oup-authoring-template'
%% with bibliographic references
%%

%%%CONTEMPORARY%%%
\documentclass[unnumsec,webpdf,contemporary,large]{oup-authoring-template}%
%\documentclass[unnumsec,webpdf,contemporary,large,numbered]{oup-authoring-template}% uncomment this line to get the numbered reference output
%\documentclass[unnumsec,webpdf,contemporary,medium]{oup-authoring-template}
%\documentclass[unnumsec,webpdf,contemporary,mediumone]{oup-authoring-template}
%\documentclass[unnumsec,webpdf,contemporary,small]{oup-authoring-template}

%%%MODERN%%%
%\documentclass[unnumsec,webpdf,modern,large]{oup-authoring-template}
%\documentclass[unnumsec,webpdf,modern,large,numbered]{oup-authoring-template}%   uncomment this line to get the numbered reference output
%\documentclass[unnumsec,webpdf,modern,medium]{oup-authoring-template}
%\documentclass[unnumsec,webpdf,modern,mediumone]{oup-authoring-template}
%\documentclass[unnumsec,webpdf,modern,small]{oup-authoring-template}

%%%TRADITIONAL%%%
%\documentclass[unnumsec,webpdf,traditional,large]{oup-authoring-template}
%\documentclass[unnumsec,webpdf,traditional,large,numbered]{oup-authoring-template}%   uncomment this line to get the numbered reference output
%\documentclass[unnumsec,webpdf,traditional,medium]{oup-authoring-template}
%\documentclass[unnumsec,webpdf,traditional,mediumone]{oup-authoring-template}
%\documentclass[namedate,webpdf,traditional,small]{oup-authoring-template}

%\onecolumn % for one column layouts

%\usepackage{showframe}
% \usepackage[backend=biber,sorting=none]{biblatex}
\usepackage{natbib}
\usepackage{amsmath}
\usepackage{algorithm}
\usepackage{algcompatible}
\usepackage{mathtools}
% \DeclareCiteCommand{~\citep}[\mkbibsuperscript]
%   {\usebibmacro{cite:init}%
%    \usebibmacro{prenote}}
%   {\usebibmacro{citeindex}%
%    \usebibmacro{cite}}
%   {\multicitedelim}
%   {\usebibmacro{postnote}}

% \addbibresource{reference.bib} 
\graphicspath{{Fig/}}

% line numbers
%\usepackage[mathlines, switch]{lineno}
%\usepackage[right]{lineno}

\theoremstyle{thmstyleone}%
\newtheorem{theorem}{Theorem}%  meant for continuous numbers
%%\newtheorem{theorem}{Theorem}[section]% meant for sectionwise numbers
%% optional argument [theorem] produces theorem numbering sequence instead of independent numbers for Proposition
\newtheorem{proposition}[theorem]{Proposition}%
%%\newtheorem{proposition}{Proposition}% to get separate numbers for theorem and proposition etc.
\theoremstyle{thmstyletwo}%
\newtheorem{example}{Example}%
\newtheorem{remark}{Remark}%
\theoremstyle{thmstylethree}%
\newtheorem{definition}{Definition}

%\usepackage{draftwatermark}
%\SetWatermarkText{Draft}
%%Uncomment  \usepackage{draftwatermark} and \SetWatermarkText{Draft} to watermark your PDF as a draft
\begin{document}

\journaltitle{Journal Title Here}
\DOI{DOI added during production}
\copyrightyear{YEAR}
\pubyear{YEAR}
\vol{XX}
\issue{x}
\access{Published: Date added during production}
\appnotes{Paper}

\firstpage{1}

%\subtitle{Subject Section}

\title[Short Article Title]{Multilateration-Based Indexing and Navigation for Error-Tolerant Read Mapping}

\author[1,$\dagger$]{Luting Zhou}
\author[1,$\dagger$]{Minghao Fang}
\author[2]{Yiru He}
\author[1]{Cheng Wang}
\author[1]{Jinpu Cai}
\author[1]{Zhenwei Huang}
\author[3]{Yuxuan Wang}
\author[4]{Shiyang Yu}
\author[5,6]{Shiyang Ma}
\author[7]{Yan Li}
\author[3,$\ast$]{Hongyi Xin}

\address[1]{\orgname{Global College, Shanghai Jiao Tong University}, \orgaddress{\city{Shanghai}, \country{China}}}

\address[2]{\orgname{School of Medicine, Shanghai Jiao Tong University}, \orgaddress{Shanghai}, \country{China}}

\address[3]{\orgname{Global Institute of Future Technology, Shanghai Jiao Tong University}, \orgaddress{\city{Shanghai}, \country{China}}}

\address[4]{\orgname{Shanghai Sixth People's Hospital Affiliated to Shanghai Jiao Tong University School of Medicine}, \orgaddress{\city{Shanghai}, \country{China}}}

\address[5]{\orgname{Institute of Clinical Medicine, Shanghai Jiao Tong University School of Medicine}, \orgaddress{\city{Shanghai}, \country{China}}}

\address[6]{\orgname{School of Mathematical Sciences, Shanghai Jiao Tong University}, \orgaddress{\city{Shanghai}, \country{China}}}

\address[7]{\orgname{State Key Laboratory of Genome and Multi-omics Technologies, BGI Research}, \orgaddress{\city{Shenzhen}, \country{China}}}




\corresp[$\ast$]{Corresponding author. \href{mailto:hongyi.xin@sjtu.edu.cn}{hongyi.xin@sjtu.edu.cn}} 



\received{Date}{0}{Year}
\revised{Date}{0}{Year}
\accepted{Date}{0}{Year}

%\editor{Associate Editor: Name}

%\abstract{
%\textbf{Motivation:} .\\
%\textbf{Results:} .\\
%\textbf{Availability:} .\\
%\textbf{Contact:} \href{name@email.com}{name@email.com}\\
%\textbf{Supplementary information:} Supplementary data are available at \textit{Journal Name}
%online.}

\abstract{Indexing is a fundamental step in read mapping.
To maximize efficiency, ideal indexers must utilize long seeds while remaining robust to sequencing errors and genetic variations.
Existing schemes typically rely on probabilistic error tolerance at fixed sub-locations or approximate edit distances via Euclidean embeddings.
However, because edit-distance space is a metric space but not a normed space, approximations using inner-product spaces (like Euclidean space) are inherently lossy.
Here, we demonstrate that high-precision positioning is attainable in coordinate-less string space through multilateration, provided the discrete space is covered by a meticulously selected set of beacon strings.
We show that for any query, a comprehensive set of high-similarity candidates can be efficiently localized using a small subset of proximal beacons.
Leveraging these insights, we propose NavigaMer, a multi-tiered navigation system for string indexing.
NavigaMer retrieves candidate reference sequences by progressively narrowing the search space within the sequence universe, guided by beacon sets ranging from coarse-to-fine resolutions.
Experiments on simulated and real genomes demonstrate that NavigaMer is fully error-tolerant and highly accurate.
It identifies all reference candidates within a specified edit distance with minimal false positives, providing a high-sensitivity foundation for future read mappers.} 

\keywords{Read mapping, fuzzy string indexing, multilateration, edit distance, seed and extend}
% \keywords[Abbreviations]{abbreviation1, abbreviation2, abbreviation3, abbreviation4}

% \otherabstract[Additional Abstract]{Use this element for elements such as Graphical abstract, Lay summary, Translated abstract etc. Que cum aut etum qui ium dolupta ssequia autati odis demporepe ad et es alit rem repudaerae min et volorum re volupta nobit volectur aut fuga.}

% \otherabstract[Graphical Abstract]{\colorbox{black!20}{\hbox to 0.97\textwidth{\vbox to 50pt{}}}}

% \boxedtext{Key Messages}{
% \begin{itemize}
% \item Key boxed text here.
% \item Key boxed text here.
% \item Key boxed text here.
% \end{itemize}}

\maketitle

%\begin{epigraph}
%Epigraph text. Ximporem qui reperov idempedit modio. Bisto imagnatem quae aceptis
%nobitae quid eum rae adignis quias-sit vellacc uptatur sunt quis rentis eaquasit alia deliquam
%rec-to consed unt. Empor sum ratur ressimusdae. Nam fugiae.
%\source{Epigraph source}
%\end{epigraph}


\section{Introduction}

The field of bioinformatics has been transformed by the advent of high-throughput sequencing,
which allows researchers to decode the genetic blueprints of organisms at unprecedented speeds. 
At the heart of this genomic revolution lies DNA read mapping,
the computational process of aligning millions of short DNA fragments (reads) to a known reference genome.
This task is foundational for critical applications, including population genetics, where it helps track evolutionary patterns, and cancer research,
where it identifies the subtle mutations driving tumor growth~\citep{icgc2020pan, li2018minimap2, zook2020robust}.
The ultimate goal of sequence alignment is to identify the specific location in the reference genome that is most similar to a query read.
However, the identified reference sequence is frequently imperfect due to sequencing errors and genetic heterogeneity.
However, the identified reference sequence is frequently imperfect due to sequencing errors and genetic heterogeneity~\citep{li2011statistical}.
Thus, A robust mapper must account for minor disparities caused by both inherent genetic variations and technical sequencing errors~\citep{ibrahim2024bioinspired}.

To handle the sheer scale of modern reference genomes, most mappers employ the seed-and-extend paradigm.
Instead of performing a computationally expensive global alignment for every read,
mappers extract shorter substrings, or seeds, to quickly identify potential candidate regions.
This strategy relies on the Pigeonhole Principle to maintain error tolerance~\citep{navarro2001guided}.
Formally, if a read of length $L$ is partitioned into $n$ non-overlapping seeds,
and we wish to tolerate up to $k$ errors, the mapper must ensure that at least one seed remains error-free.
This is guaranteed if: $n > k$.
By selecting $k+1$ seeds, the mapper ensures that even with $k$ substitutions or indels, at least one seed will result in an exact match (a ``hit'') within the reference index,
triggering the more intensive extension phase to verify the full alignment~\citep{xin2016optimal}.


\begin{figure*}[t]
    \centering
    \includegraphics[width=\textwidth]{oup-authoring-template/Figure1.jpg}
	\caption{Overview of NavigaMer. (a) Illustration of the DNA sequence space $\mathcal{M}$ and edit distance. (b) The query process for identifying a query sphere (i.e., finding sequences in the reference genome that are within a specified distance tolerance from the query). (c) Schematic diagram of the fuzzy multilateration process and its resulting envelope. (d) Left: The hierarchical structure of NavigaMer. Right: Beacons within each tier of the world used for local fuzzy multilateration.}
    \label{Fig:overview}
	\label{fig1}
\end{figure*}


Read mappers must navigate a delicate trade-off between indexing efficiency and error resilience.
For maximum efficiency, a mapper should utilize longer seeds ($k$-mers),
as the probability of a specific substring occurring by chance in a genome of size $G$ decreases exponentially with its length.
Longer seeds result in fewer ``spurious'' hits, significantly reducing the search space.
However, as seed length increases, the mapper’s resilience to errors inevitably drops.
Since the total length of the read is fixed, using longer seeds reduces the number of non-overlapping segments available,
making it statistically more likely that every seed will be corrupted by at least one sequencing error or variant~\citep{xin2015shifted}.

Several innovative strategies have attempted to reconcile this dilemma by allowing seeds themselves to be error-tolerant:
\begin{itemize}
\item Spaced Seeds: This approach utilizes a binary mask to ignore specific positions within a seed.
By focusing only on ``match'' positions and allowing substitutions at ``do not care'' positions, spaced seeds increase the probability of a hit in the presence of point mutations~\citep{ma2002patternhunter}~\citep{kazemi2022nthash2}.
\item Strobemers: Extending logic of spaced seeds to include insertions and deletions (indels), strobemers link multiple short $k$-mers (sub-seeds) separated by variable-length gaps~\citep{sahlin2021effective}~\citep{sahlin2022strobealign}~\citep{jain2020weighted}.
By using a ``minstrobe'' selection based on minimum hashing, this method remains flexible to indels between the sub-segments.
\item TensorSketch: This method shifts the problem from string space to vector space.
It converts seeds into high-dimensional vectors in a Euclidean space, where the Euclidean distance serves as an approximation of the Edit Distance~\citep{joudaki2023aligning}.
These vectors are then indexed using a Hierarchical Navigable Small Worlds (HNSW) data structure to find proximal neighbors~\citep{malkov2018efficient}. 
\end{itemize}

Despite these advancements, existing methods remain either partially error-tolerant or mathematically limited.
Spaced seeds and strobemers are still fragile; a single error occurring in a match position or within the strobe itself will cause the seed to fail.
On the other hand, TensorSketch uses high-dimensional embeddings, approximating the discrete edit distance between sequences with continuous Euclidean distance. 
Yet, while DNA sequences and edit distance define a strict metric space, they do not constitute a normed space, meaning any projection into a Euclidean inner-product space is mathematically forced to incur inherent lossiness and distortion~\citep{ostrovsky2007low, yuan2025melo, marccais2019locality}. 
This fundamental incompatibility means that any Euclidean approximation will inevitably produce false negatives or misrepresent the true biological distance between sequences.
Moreover, the reliance on HNSW introduces a secondary source of information loss; as an approximate nearest neighbor algorithm driven by greedy graph traversal, it inherently trades perfect recall for search efficiency, risking the omission of true biological matches if the routing gets trapped in local optima~\citep{munyampirwa2024down}.

In this paper, we propose NavigaMer, a novel indexing strategy that sketches an enclosure of candidate sequences directly within the edit-distance string space.
Instead of lossy embeddings, NavigaMer leverages the triangle inequality property of metric spaces to bound the search area ~\citep{kulis2013metric}.
By calculating the distance of a query against a meticulously pre-selected set of ``beacon'' strings,
we can mathematically prune the sequence universe~\citep{chavez2001searching, yianilos1993data, chen2017pivot}.
We demonstrate that by employing only a small subset of proximal beacons, NavigaMer achieves both high precision and perfect recall.
Based on this observation, we introduce a hierarchical multi-tiered positioning system.
This system localizes query candidates through progressive refinement,
moving from coarse-grained to fine-grained beacon sets.

% The contributions of this work are threefold. First, we propose NavigaMer, a hierarchical indexing strategy that replaces lossy embeddings with precise geometric multilateration, ensuring that the search process respects the non-Euclidean nature of the edit distance. Second, we provide a theoretical proof that our beacon-based pruning guarantees the inclusion of all valid candidates, eliminating false negatives caused by heuristic failures. Third, extensive simulations demonstrate that NavigaMer maintains 100\% recall at 15\% mutation rates——a regime where traditional seed-based aligners fail. This establishes NavigaMer as a robust foundation for next-generation error-tolerant read mapping.
%NavigaMer thus provides a high-sensitivity foundation for the next generation of read mappers, ensuring no valid biological candidate is left behind.

%%%%%%%%%%%%%%分割线%%%%%%%%%%%%%%



\section{Methods}

\subsection{The Coordinate-less Geometry of the DNA Sequence Space and Fuzzy Multilateration} 

Fundamentally, indexing can be perceived as a localization problem within a discrete, coordinate-less metric space illustrated in Figure~\ref{fig1}.\textbf{a}: 
Given a query sequence $Q$, an edit distance threshold $t$, and a large set of reference sequences $\textbf{G}$ that contains all the DNA snippets of length $|Q|$ extracted from a reference genome, 
find a subset $\textbf{C}_Q \subset \textbf{G}$ such that: $\forall R \in \textbf{G}$ where $\mathrm{Dist}_{\mathrm{edit}}(Q, R) \leq t$, we have $R \in \textbf{C}_Q$.
If the above condition is satisfied, we claim that the index retrieval $\textbf{C}_Q$ is comprehensive.
If the number of sequences $R \in \textbf{C}_Q$ such that $\mathrm{Dist}_{\mathrm{edit}}(Q, R) > t$ is small, we claim that $\textbf{C}_Q$ is tight. 
Specifically, if such $R$ does not exist, we claim that $\textbf{C}_Q$ is exact.
Note that the territory of a tight $\textbf{C}_Q$ is defined by both a geometric center $Q$ and a distance radius $t$.
Therefore, as demonstrated in Figure~\ref{fig1}.\textbf{b}, it can be perceived as a sphere from a geometric perspective.
Nonetheless, unlike canonical geometry systems such as Euclidean or Riemannian geometry, which are analytic (coordinate geometry) and encompass an infinite space,
the DNA sequence space is discrete, finite, and inherently coordinate-less~\citep{de2008computational}.
Thus, it is a synthetic geometry~\citep{bertrand2021stratified, mondino2025equivalence}.

Due to its synthetic geometric nature, range positioning in the DNA geometry is fundamentally different from the same task in analytic geometry.
In GPS or radar systems, range positioning is achieved by utilizing a set of anchors, concretely as finding an enclosure from a system of algebraic equations associated with coordinates~\citep{malivert2023comparison, zafari2019survey}.
Such a practice is not applicable to the DNA sequence space, which is not defined by coordinates.
However, in the special case of measuring sequence similarity with edit distances, the DNA sequence space still satisfies the criteria of being a metric space~\citep{kulis2013metric}.
Hence, range positioning, which is essentially identifying the sphere $\textbf{C}_Q$, 
can still be carried out following the same philosophical principles of a multilateration GPS positioning system,
but implemented using a discrete mathematical approach.
As we will show below in this section, fuzzy multilateration in the reference sequence space $\textbf{G}$, as demonstrated in Figure~\ref{fig1}.\textbf{c}, utilizes specific sequences as anchors, which we call beacons in our terminology, to produce a comprehensive, though not necessarily tight, envelope $\textbf{E}_Q$ of $\textbf{C}_Q$, meaning that $\textbf{C}_Q \subset \textbf{E}_Q$.
% The tightness of $\textbf{E}_Q$ is closely related with the selection of multilateration anchors, which we call beacons in this manuscript.
% We then show that a relatively tight, while comprehensive $\textbf{E}_Q$ can be obtained with a set of meticulously selected beacon set $\textbf{B}_Q$ for a particular query $Q$.

Formally, we define the geometric features and operations within the metric space of DNA sequences in Figure~\ref{fig1}.\textbf{a} in \textbf{Appendix A}.

% \begin{definition}
% Let $\Sigma = \{A, T, C, G\}$ be the nucleotide alphabet, and $\mathcal{S} = \Sigma^*$ represents the set of all possible nucleotide sequences. Let $\mathrm{Dist}_{\mathrm{edit}}: \mathcal{S} \times \mathcal{S} \to \mathbb{R}_{\geq0}$ be the edit distance between two sequences. We define the global DNA sequence space $\mathcal{M}$ as a metric space:
% \begin{align}
%     \mathcal{M} \coloneqq (\mathcal{S}, \mathrm{Dist}_{\mathrm{edit}}).
% \end{align}
% Within $\mathcal{M}$, the reference genome pool $\textbf{G} \subset \mathcal{S}$ is a finite subset. The exact target candidate set, or the \textit{query sphere }$\textbf{C}_Q$ for a query $Q \in \mathcal{S}$ under error threshold $t \in \mathbb{R}_{\geq0}$ is equivalent to the intersection of $\textbf{G}$ and the closed ball centered at $Q$:
% \begin{align}
%     \textbf{C}_Q \coloneqq \{ R \in \textbf{G} \mid \mathrm{Dist}_{\mathrm{edit}}(R, Q) \leq t \}.
% \end{align}
% \end{definition}

% % In metric space $\mathcal{M}$, the envelope $\textbf{E}_Q$ is determined by the inherent triangular inequality against a set of selected beacons.
% In metric space $\mathcal{M}$, each beacon $b \in \mathcal{S}$ could provide a specific constraint set, geometrically as an annulus, that contains $\textbf{C}_Q$.

% \begin{definition}
%     For a given beacon sequence $b \in \mathcal{M}$, a query sequence $Q \in  \mathcal{M}$, and an edit distance threshold $t \in \mathbb{R}_{\geq0}$, we define the \textit{annulus} $\textbf{A}_{b, t}(Q)$ as the set of sequences in $\mathcal{M}$ whose distance to $b$ is within $t$ of the distance between $Q$ and $b$:
%     \begin{align}
%        \textbf{A}_{b, t}(Q) \coloneqq \{ u \in \mathcal{M} \mid |\mathrm{Dist}_{\mathrm{edit}}(Q, b) - \mathrm{Dist}_{\mathrm{edit}}(u, b)| \leq t \}.
%     \end{align}
%     \label{def:annulus}
% \end{definition}

% \begin{theorem}
% For any arbitrary beacon $b \in \mathcal{M}$, query $Q \in \mathcal{M}$, and threshold $t \in \mathbb{R}_{\geq0}$, then the query sphere is bounded by the annulus:
% \begin{equation}
%     \textbf{C}_Q \subseteq \textbf{A}_{b, t}(Q) \cap \textbf{G}.
% \end{equation}
%     \label{thm:ballinann}
% \end{theorem}


% Building upon Theorem \ref{thm:ballinann}, fuzzy multilateration is essentially determining the envelope $\textbf{E}_Q$ through the geometric intersection of annuli derived from multiple beacons. Thereby, $\textbf{C}_Q \subseteq \textbf{E}_Q$.

% \begin{definition}
% Let $\textbf{B}(Q) \subset \textbf{G}$ be a finite set of reference sequences serving as beacons for query $Q$. The fuzzy multilateration produces \textit{envelope} $\textbf{E}_{Q, \textbf{B}(Q)}$, or briefly $\textbf{E}_Q$, as:
% \begin{equation}
%     \textbf{E}_{Q, \textbf{B}(Q)} \coloneqq \textbf{G} \cap \bigcap_{b \in \textbf{B}(Q)} \textbf{A}_{b, t}(Q).
% \end{equation}
% \label{def:multilateration}
% \end{definition}

% \begin{remark}
%     For any query $Q$ and beacon set $\textbf{B}(Q)$, the envelope $\textbf{E}_Q$ strictly preserves the comprehensive target candidate set: $\textbf{C}_Q \subseteq \textbf{E}_Q$.
%   \label{cor:includeall}
% \end{remark}
% \begin{proof}
% Since Theorem \ref{thm:ballinann} holds for every individual beacon $b \in \textbf{B}(Q)$, it must hold that $\textbf{C}_Q \subseteq \textbf{G} \cap \bigcap_{b \in \textbf{B}(Q)} \textbf{A}_{b, t}(Q)$.
% \end{proof}



\subsection{Maximum Beacon Coverage and Close-Range Beacon Positioning}

%%%%%%%%%%%%%%辛%%%%%%%%%%%%%%

Unlike algebraic multilateration in $d$-dimensional Euclidean geometry,
where accurate positioning is guaranteed given distances to any $d+1$ unique anchors with known coordinates,
positioning through multilateration within the discrete DNA sequence space is rarely exact due to its synthetic geometry properties ~\citep{malivert2023comparison, zafari2019survey, bertrand2021stratified}.
In this metric space, positioning efficiency is inextricably linked to both the quantity and the proximity of beacons.
The tightness of the multilateration envelope $\mathbf{E}_Q$ in approximating the true candidate set $\mathbf{C}_Q$ depends heavily on the beacon selection $\mathbf{B}(Q)$.
In an ideal case, where a query $Q$ is already present in the reference genome $\mathbf{G}$ and is itself selected as a beacon ($Q \in \mathbf{G}$ and $Q \in \mathbf{B}(Q)$),
the envelope becomes exact: $\mathbf{E}_Q = \mathbf{C}_Q$.
More generally, as the number of beacons $|\mathbf{B}(Q)|$ increases, the enclosure asymptotically approaches the candidate set ($\mathbf{E}_Q \rightarrow \mathbf{C}_Q$).
The placement of these beacons is equally consequential and highly query-dependent.
Specifically, the enclosure $\mathbf{E}_Q$ grows tighter as the beacons in $\mathbf{B}(Q)$ are drawn from the query’s immediate vicinity;
that is, $\mathbf{E}_Q \rightarrow \mathbf{C}_Q$ as $\mathrm{Dist}_{\mathrm{edit}}(b, Q)$ decreases for all $b \in \mathbf{B}(Q)$.
These observations suggest that for cost-effective multilateration, where the goal is to draft a sufficiently tight enclosure $\mathbf{E}_Q$ using a minimal number of edit-distance computations,
the beacon set $\mathbf{B}(Q)$ must be dynamically selected and tailored to the unique local sequence space topological profile of each individual query.
Specifically, for each query $Q$, the optimal strategy is to select a few proximal beacons to form the beacon set $\mathbf{B}(Q)$.

To demonstrate the above, we formalize the definition of tightness and mathematically validate the impacts of beacon quantity and proximity to determining $\mathbf{C}_Q$ on the envelope $\textbf{E}_Q$ in \textbf{Appendix B}.


\subsection{NavigaMer Overview}

As demonstrated in in Figure~\ref{fig1}.\textbf{d}, NavigaMer is essentially a multi-tiered multilateration scheme, where it iteratively and progressively tightens the candidate selection envelope $\textbf{E}_Q$ of a query $Q$.
To achieve this, NavigaMer partitions the reference sequence universe $\textbf{G}$ into overlapping pockets of ``worlds'',
where each world $\textbf{W}_C$ is a local proximal and continuous spherical subspace of $\textbf{G}$ defined by a geometric center $C \in \textbf{G}$ and a radius $r$.
The worlds are organized into separate tiers, from coarse grain to fine grain, ranked by their radii.
This world system is strictly inclusive, meaning that any smaller world must be fully included by an immediate larger world, forming a parent-children hierarchical relationship.

Borrowing an analogy from astronomy, this world hierarchy can be perceived as galaxies, systems, planets and continents.
%, borrowing the analogy from astronomy.
Each world $\textbf{W}_C$ contains a set of beacons $\textbf{B}_{\textbf{W}_C}$ to support within-world navigation via multilateration.
A lookup, therefore, can be perceived as a progressive navigation within this world hierarchy:
Given a query $Q$ and a larger world $\textbf{W}_L$ that fully contains $\textbf{C}_Q$,
NavigaMer quickly identifies the immediate sub-world $\textbf{W}_S \subset \textbf{W}_L$ that also fully engulfs $\textbf{C}_Q$ ($\textbf{C}_Q \subset \textbf{W}_S$), and enters $\textbf{W}_S$ for more fine-grained navigation.
As NavigaMer travels deeper into the world hierarchy, it also gradually confines the multilateration to smaller proximities of the query $Q$,
using beacons that are progressively closer to $Q$.
In this manner, NavigaMer confines the scope of localization and limits the scale of multilateration workload in each step,
while quickly homing the navigation into a small proximal terminal pocket of $Q$ in $\textbf{G}$.
In the terminal small world, only a small set of close-distance beacons are employed for a final-round of multilateration against a limited number of reference sequences in this terminal world.

%In each iteration, NavigaMer 


% -----
% 目标：全局方法感知。空间长什么样？
% unique strings in a referrence; into a multi-tier tree-like space
% query 是 trasverse through hierarchy
%  我们只选 proximal beacon 进行 multi
%  分割；by exploiting metric space property, world into pockets
% narrow down search space
% 层次性质
% 世界中的 local beacon（不同 resolution）
% 多层  multilateration


\subsection{Multi-Tiered Beacon Navigation via Hierarchical Multilateration}

% 讲怎么multilateration————一个世界里到底存了什么（比如小世界的中心？），然后multilateration的具体算法。具体！具体！具体！

% 讲如果不是fully inclusive的话怎么做。

NavigaMer models the query process as a top-down hierarchical multilateration. Rather than exhaustively computing sequence alignments, the algorithm progressively confines the search space by utilizing the pre-computed spatial topography of the world graph to form a tight envelope $\textbf{E}_Q$ that approximates the exact query sphere $\textbf{C}_Q$.
% 人话: 不是fully inclusive怎么做
NavigaMer executes a top-down hierarchical multilateration by structuring each index "world" to store its center, radius, local beacons, and the exact metric bounding boxes (MBBs) of its sub-worlds. During a query, the algorithm utilizes $O(1)$ distance checks against these beacons to instantly prune any sub-world whose MBB, even when expanded by the error tolerance, cannot mathematically intersect the query sphere. Crucially, if the query sphere is not fully contained within a single sub-world but instead straddles multiple boundaries, the search dynamically branches to evaluate all intersecting sub-worlds, effectively preventing false negatives and guaranteeing comprehensive retrieval.

Technical description and details of our query process are provided in \textbf{Appendix C}.

% \begin{algorithm}[t]
% \caption{Hierarchical Multilateration Search}
% \label{alg:navigation}
% \begin{algorithmic}[1]
% \REQUIRE Query $q$, Tolerance $\varepsilon$, Index Root $Root$
% \ENSURE Matching sequences $Result$
% \STATE $W_{current} \leftarrow \{Root\}$
% \STATE /* Phase 1: Top-Down Pruning */
% \FOR{$l = 1$ \textbf{to} $k-1$}
%     \STATE $W_{next} \leftarrow \emptyset$; \quad $FastPath \leftarrow \text{FALSE}$
%     \FORALL{$w \in W_{current}$}
%         \FORALL{Child $c$ of $w$}
%             \STATE \textbf{/* Prune using pre-computed MBBs */}
%             \IF{$\exists b \in c.\text{Beacons}$ s.t. $\text{dist}(q, b) \notin [c.\text{MBB}_b.\text{min} - \varepsilon, c.\text{MBB}_b.\text{max} + \varepsilon]$}
%                 \STATE \textbf{Continue}
%             \ENDIF
%             \STATE $dist \leftarrow d(q, c.\text{center})$
%             \IF{$dist \le c.R + \varepsilon$} 
%                 \STATE $W_{next} \leftarrow W_{next} \cup \{c\}$
%                 \IF{$dist + \varepsilon \le c.R$} \COMMENT{Strict Containment}
%                     \STATE $W_{next} \leftarrow \{c\}$; $FastPath \leftarrow \text{TRUE}$; \textbf{Break}
%                 \ENDIF
%             \ENDIF
%         \ENDFOR
%         \IF{$FastPath$} \STATE \textbf{Break} \ENDIF
%     \ENDFOR
%     \STATE $W_{current} \leftarrow W_{next}$
% \ENDFOR

% \STATE \textbf{/* Phase 2: Leaf Refinement */}
% \FORALL{Leaf $w \in W_{current}$}
%     \FORALL{$s \in w.\text{Data}$}
%         \IF{$\forall b \in w.\text{Beacons}, |d(q, b) - d(s, b)| \le \varepsilon$}
%             \IF{$d(q, s) \le \varepsilon$} \STATE $Result \leftarrow Result \cup \{s\}$ \ENDIF
%         \ENDIF
%     \ENDFOR
% \ENDFOR
% \end{algorithmic}
% \end{algorithm}

\subsection{Beacon Set Selection and Hierarchical Sequence Universe Topography}

NavigaMer builds the world hierarchy while simultaneously selects multilateration beacons at the same time through construction by insertion.
Briefly, during the world hierarchy construction, NavigaMer simultaneously strives towards two goals:
1) All worlds are well-sprinkled with and covered by beacons, such that for any sub-world, there is always a close-by beacon in its vicinity.
2) The world hierarchy is perfectly maintained, such that a parent world is properly associated with all next-tier sub-worlds that either partially or fully overlaps with the parent world.

To ensure these properties,
NavigaMer builds the world hierarchy in 3 stages.
In the first stage,
NavigaMer constructs an extended world hierarchy sketch that has a tier count twice of the world tiers,
with a temporary immediate tier between any two adjecent world tiers.
New worlds across tiers are created if an insertion (and its associated new world centered around it) is not fully covered by an existing world (or super-world).
In this stage, the comprehensive parent-children relationship is not guaranteed yet.
Instead, a sub-world is only connected to one of its subsuming super-world.
In the second stage, an additional round of reference sequence insertion is repeated.
In this round, no new world is created.
Instead, it reviews all parent and children relationships and reconnects the missing links,
producing a graph rather than a tree,
where a super-world is connected with all sub-worlds that partially or fully overlaps with it.
In the final stage,
the world centers of the temporary intermediate tiers between world tiers are extracted as multilateration beacons for the parent world.
The intermediate tiers are then removed while repairing the parent-children relationships.

By extracting the intermediate tier as multilateration beacons,
NavigaMer ensures uniform spreads of beacons in worlds spanning across the entire sequence world.
It also guarantees that there exist at least one beacon in an immediate vicinity for any reference sequence or sub-world.
The technical details of world construction is demonstrated in \textbf{Appendix D}.

\section{Results}

\begin{figure}[t] % 移除星号，变为单栏环境
    \centering
    \includegraphics[width=\columnwidth]{oup-authoring-template/Figure2.png} % 宽度改为 \columnwidth 以适配单栏
    \caption{NavigaMer benchmark with other methods. (a) Recall changing according to mutation count for NavigaMer, Stromeber, Spaced Seed and TensorSketch. (b) Precision-Recall tradeoff curve for different methods. (c) Candidate quality distribution and candidate distance for different methods.}
    \label{fig2}
\end{figure}

\subsection{Benchmark Setup}
We implemented the NavigaMer index and search algorithms in C++ All experiments were conducted on a Linux server (Kernel 6.17) running Python 3.11. The benchmarking environment utilized an Intel CPU with 96 logical cores. For benchmarking, we consider two architectural families: seed-and-extend methods (Strobemers, Spaced Seeds) that decompose a query read into short overlapping seeds and report every genomic locus at which a seed matches, and sketch-based methods TensorSketch that embed the full read into a compact sketch and retrieve the nearest neighbours from an approximate nearest-neighbour (ANN) index.

\subsubsection{Datasets and Simulation}
To evaluate the indexers on a highly complex and repetitive sequence space, we utilized the human reference genome (GRCh38) as our reference dataset. Synthetic query reads (100 bp) were generated by randomly sampling subsequences across the genome. To comprehensively simulate natural genetic variations and sequencing errors, mutations (encompassing substitutions, insertions, and deletions) were introduced at varying edit distances. All dataset generations and mutational events were executed using a deterministic random seed to guarantee strict reproducibility of the ground truth multi-mapping coordinates.

\subsubsection{Evaluation Metrics and Protocol}
To rigorously evaluate the comprehensiveness and filtration efficiency of the indexers under our geometric framework, we establish locus-level evaluation metrics. For a given query $Q$, let the ground truth $\textbf{C}_Q$ be the exact target candidate set, ensuring locus-level comprehensiveness for multi-mapped queries. Let $\textbf{O}_Q$ be the candidate set retrieved by an arbitrary indexer (where $\textbf{O}_Q = \textbf{E}_Q$ for NavigaMer). A retrieved candidate $R' \in \textbf{O}_Q$ successfully captures a truth sequence $R \in \textbf{C}_Q$ if their physical genomic coordinates $\mathrm{loc}(\cdot)$ are within a tolerance threshold $\delta$, such that $|\mathrm{loc}(R) - \mathrm{loc}(R')| \leq \delta$.

To quantify comprehensiveness, we measure the indexer's Recall, defined as the fraction of truth loci in $\textbf{C}_Q$ successfully captured by at least one candidate in $\textbf{O}_Q$. Any truth locus failing this hit condition is classified as a False Negative ($FN$). Furthermore, to assess filtration efficiency (tightness), we eschew traditional precision metrics, which can be misleading for sequence indexers. Instead, we directly measure the overall candidate volume $|\textbf{O}_Q|$. Candidates in $\textbf{O}_Q$ that fail to capture any sequence in $\textbf{C}_Q$ constitute False Positives ($FP$). A geometrically tighter indexer produces a significantly smaller candidate volume $|\textbf{O}_Q|$ while maintaining high Recall, effectively minimizing the $FP$ overhead and drastically reducing the computational burden for downstream dynamic programming verification.


% 错误多，容错概率降；我们不降。
% Precision不够好（经常给出垃圾-edit 大）
% Recall差（不是所有candidate都找到）
% 怎么都调不好（Precision vs. Recall）

\subsection{Benchmark Results}

\subsubsection{Comprehensiveness Under High Mutation Rates}
NavigaMer is algorithmically designed to ensure comprehensiveness. To evaluate robustness, we assessed locus-level recall across mutation counts ranging from 1 to 15 under fixed parameters. Every valid genomic locus satisfying the edit distance threshold $\epsilon$ constitutes an independent ground truth. 

As shown in our results Figure~\ref{fig2}.\textbf{a}, NavigaMer's geometric indexing provides a theoretical guarantee for locus retrieval, empirically maintaining a perfect recall of $1.0$ across the entire tested spectrum. In contrast, all baseline methods exhibit significant performance degradation as sequence divergence increases. Strobemer shows the sharpest decline due to the failure of consecutive seed matching in hyper-variable regions. While Spaced Seeds and Tensor Sketch perform moderately better, they intrinsically fail to sustain comprehensive retrieval, leaving substantial false negatives under high error conditions.

\subsubsection{Precision-Recall Trade-off}
Navigamer does not sacrifice precision for recall, since it provides a close envelope over the query sphere. A fundamental limitation of existing methods is that they achieve high recall strictly at the expense of precision, heavily inflating the candidate volume. To systematically evaluate this trade-off, we swept the key hyperparameters governing match stringency and candidate volume for all methods (e.g., seed length/weight, window offsets, and top-$k$ retrieval). 

As shown in Figure~\ref{fig2}.\textbf{b}, for seed-based methods, relaxing parameters (e.g., reducing strobe length $k$ or seed weight $\omega$) drives recall upward but floods the candidate set with false positives, causing a severe collapse in precision. TensorSketch exhibits a similar bottleneck: expanding the neighbor budget ($k_{\mathrm{ANN}}$) steadily degrades precision before recall can saturate. Conversely, NavigaMer fundamentally breaks this trade-off. By utilizing exact geometric filtration bounded by the edit-distance tolerance, NavigaMer operates in a uniquely optimal regime——guaranteeing near-perfect recall while strictly preserving high precision and maintaining a compact candidate set.

\subsubsection{Candidate Quality Guarantee}
NavigaMer provides a high quality candidate set, returning all close-up sequences. To further demonstrate that our method preserves candidate quality, we analyzed the Levenshtein edit distance distribution between query reads and retrieved candidates when all indexers are configured to maximize recall.

NavigaMer guarantees that candidate distances tightly mirror the true mutation profile. As observed in Figure~\ref{fig2}.\textbf{c}, the edit distances of NavigaMer's candidates are strictly bounded near the query ($\leq 15$ edits), demonstrating a highly tight retrieval envelope. In stark contrast, baseline methods exhibit severe precision loss. Seed-based heuristics yield pronounced "long-tail" distributions burdened with distant false positives, while TensorSketch produces a massive secondary density peak at extreme edit distances (50--60), exposing the inaccuracies of its spatial embedding. Ultimately, NavigaMer's superior filtration guarantees the retrieval of high-quality, proximal candidates, directly minimizing the computational overhead for downstream dynamic programming verification.
% \textbf{Uncompromised Robustness at High Mutation Rates}
% % 关注容错率————随着错误率变高，不同方法的容错效果怎么变？
% % 为了达到容错，不同方法牺牲的效果如何？
% \textbf{Comprehensiveness.}
% % 方法是否能找到每个 query 的至少一个匹配？（只要有一个匹配，read 就算一个 FP）
% % 方法是否能找到每个 query 的所有匹配？（每个匹配算一个 TP）
% \textbf{Candidate quality.}
% % 找到的序列质量如何？（和query 的相似度如何？）




% 验证：beacon多，envolope好。
% beacon近，envolope好。
% 少数近beacon，几乎就最好。

% 空间分割法：
% Overlap情况。
% 要split search的情况。
% 各种不同建图参数（世界的半径，beacon密度）对于split map……等等
% Corner case情景。(optional)
% 逐渐扩大，调控search space reduction率。

% 对于reference的meta分析。

% 速度。

% 实验目标(@yiru)
% 什么样的 beacon 能够最快缩小 search space？beacon 的选择怎样更好？

% 实验设计(@yiru)
% 

% 实验结果及结论(@yiru)
% 越多越好；越近越好；少数近的跟很多远的效果相近

% To assess the robustness of our method under varying degrees of genetic divergence, we simulated the search space reduction process at mutation rates of 5\%, 15\%, and 25\% using a $10,000$ bp random reference genome. We compared a random strategy against a proximal strategy, averaging the results over $96$ trials for each condition. As shown in, the proximal strategy consistently exhibits an exponential reduction in the candidate envelope size, converging to the target region significantly faster than random selection regardless of the error density.

\subsection{Internal Evaluation and Parameter Optimization}
In this section, we demonstrate how NavigaMer's internal geometric mechanisms---specifically beacon selection, hierarchical depth, and spatial scaling---optimally reduce search space and minimize computational overhead.

% kkk

\begin{figure}[t] % 移除星号，变为单栏环境
    \centering
    \includegraphics[width=\columnwidth]{oup-authoring-template/Figure3.jpg} % 宽度改为 \columnwidth 以适配单栏
    \caption{NavigaMer internal evaluation. (a) Envelope size related to beacon selection, a close envelope can be guaranteed through several close-up beacons. (b) Impact of layer count on distance calculations, three layers is sufficient enough. (c) World access count related to world radius. (d)Access type related to tolerance, most queries will fall into a small world as a fast query.}
    \label{fig3}
\end{figure}

\subsubsection{Efficiency of Proximal Beacon Selection}
NavigaMer achieves rapid search space reduction by prioritizing proximal beacons, which inherently provide significantly stronger geometric filtering power than random sampling. We evaluated how beacon selection strategies impact the envelope size (the volume of candidate sequences satisfying the triangle inequality). By varying the beacon count, we compared a random uniform sampling strategy against a proximal strategy that strictly selects beacons with the minimum edit distance to the query. As demonstrated in Figure~\ref{fig3}.\textbf{a}, the proximal strategy drastically outperforms random selection. Notably, a minimal set of nearby beacons achieves a pruning efficiency that would otherwise require hundreds of randomly selected beacons, proving that spatial relevance is paramount for accelerating exact geometric filtration.

\subsubsection{Impact of Hierarchical Layer Count}
Structuring the spatial index into a multi-layer hierarchy drastically reduces the computational overhead of exact alignment, cutting mean distance calculations per query by nearly an order of magnitude. We analyzed the distance calculations required across different index depths. As illustrated in Figure~\ref{fig3}.\textbf{b}, transitioning from a flat, single-layer index to a three-layer hierarchy precipitously drops the calculations. The addition of a fourth layer yields a marginal further reduction. This confirms that a concise three-layer geometric architecture is sufficient to aggressively prune the search space, bypassing exhaustive evaluations without introducing unnecessary structural complexity.

\subsubsection{Optimization of World Radius Ratios}
NavigaMer minimizes index traversal costs by utilizing an optimal geometric scaling ratio between consecutive hierarchical layers, effectively balancing upper-layer encapsulation with lower-layer pruning. We investigated how the radius ratio of the middle world to the lower world ($r_{mw} / r_{sw}$) influences the mean node access counts (Figure~\ref{fig3}.\textbf{c}). Expanding the middle world radius geometrically increases the upper layer's capacity to encapsulate queries, leading to a steady decline in upper-layer accesses. However, an overly expansive middle world inherently encloses an excessive number of lower worlds, causing lower-layer accesses to surge exponentially when the ratio exceeds 6. The resulting U-shaped access distribution reveals a strict optimal scaling ratio (between 5 and 6) that minimizes the overall traversal bottleneck.

\subsubsection{Query Encapsulation Efficiency}
NavigaMer's spatial topology maximizes query encapsulation within single geometric boundaries, directly circumventing the costly overhead associated with traversing overlapping regions. We classified index accesses into two topological patterns: fully contained within a single "world" or overlapping across multiple worlds. As shown in Figure~\ref{fig3}.\textbf{d}, as the ratio of the world radius to the query edit-distance tolerance ($r_{sw} / \tau$) increases, the proportion of fully encapsulated queries grows monotonically, exceeding 80\% at a ratio of 2.5. This structural advantage guarantees that for practical biological error tolerances, the query envelope remains intrinsically localized, ensuring highly efficient retrieval.
%数量

%距离

% \subsection{Search Space Reduction Efficiency is Related to the Hierarchical Universe Topology}

% \textbf{层数}
% 会split off的频率。

% \textbf{半径}

% \subsection{Sequence Space Topology of the XXX Genome}

% 有没有重复序列、高相似序列等等。

% \subsection{Indexing Efficiency Benchmarks of Real Data}

% 和其他方法比。

% \section{Conclusion and Discussion}




\section{Discussion and Conclusion}

% While NavigaMer demonstrates exceptional efficiency and precision in mapping low-divergence, continuous short reads, we acknowledge its inherent limitations when processing complex sequence variations. For localized variations, such as short insertions or deletions (e.g., 10-30 bp), our fixed-window approach requires a correspondingly higher error tolerance (t) to rescue alignments. However, this introduces a fundamental trade-off: elevating t improves sensitivity but inevitably expands the search space, increasing computational overhead and the risk of false positives, particularly in highly repetitive genomic regions.

% Furthermore, massive structural variants, such as large transposon insertions, expose a boundary condition of our current design. Because NavigaMer relies on contiguous 150bp comparisons, a massive insertion effectively breaks the alignment window. In such scenarios, a simple adjustment of t is algorithmically insufficient to bridge the gap, distinguishing our continuous-window method from gapped-seed approaches like strobemers.

% Rather than forcing a single algorithmic step to resolve all scales of variation, we envision deploying NavigaMer as an ultra-fast, ``first-pass'' mapping engine within comprehensive analytical pipelines. Operating strictly at a low t, NavigaMer can rapidly anchor the vast majority of standard reads with minimal computational cost. The remaining unmapped fraction——potentially harboring complex structural variants——can then be strategically deferred to highly sensitive, split-read or multi-anchoring modules. This two-tiered architecture maximizes overall throughput without compromising the rigorous discovery of complex biological phenomena.
% \section{Conclusion}

In this work, we introduced NavigaMer,
a multilateration-based, error-tolerant DNA sequence indexing framework.
NavigaMer exploits the metric space property of the edit-distance space and converts the string indexing problem into a localization problem in a synthetic geometry.
By partitioning the reference genome sequences into overlapping worlds of diverse radii while connecting the worlds into a hierarchy,
NavigaMer supports cost-effective, high-accuracy fuzzy range localization in the DNA sequence space by gradually tightening a query envelop via navigating the world hierarchy through tiered multilateration.
In this manner, NavigaMer guarantees comprehensive and successful close-distance reference sequence retrieval for any query sequence that is a few-error variant any reference sequence,
regardless of the error type or error locations, as shown in Figure~\ref{fig2}.\textbf{c}.

In comparison to string-based indexing scheme such as BWT and FM-indexing,
a major downside of NavigaMer is its high volume of edit-distance computation.
At each world tier, NavigaMer has to compute the edit-distances between the query $Q$ against all beacons in that world.
Such computation is repeated in all visited worlds as NavigaMer traverses the world hierarchy.
While definitely computationally more expensive,
however, with proper cache-reuse optimizations, which we intend to explore in our future work,
NavigaMer is not necessarily associated with lower indexing throughput.
For BWT and FM-indexing,
a query involves $O(|Q|)$ random memory accesses with little-to-none data reuse.
In contrast, in NavigaMer,
all query sequences that access the same world uses the same beacon set.
With proper scheduling when processing bulk indexing operations by grouping queries to the same world together,
NavigaMer could potentially greatly benefit from high degree of data reuse.
Considering that read mapping is already memory-latency-and-throughput bound~\citep{GenASM, onurframework},
the cache-friendliness of NavigaMer could potentially translate to higher indexing throughput than BWT and FM-indexing.
We plan to further optimize the bulk lookup procedure of NavigaMer in our future work.

Another potential limitation of NavigaMer is its fixed length indexing scheme.
Unlike BWT and FM-indexing which supports lookup for arbitrary length of $Q$,
NavigaMer only supports lookups for fixed full-length query sequences.
While this is generally acceptable for query sequences that contains only single nucleotide variations,
NavigaMer in its default setting is incompatible with split-mapping for RNA transcripts or structure variations.
One potential remedy for split mapping is to employ multiple NavigaMer instances at various lengths:
if a sequence cannot find any proper small-distance candidates,
indicating that the sequence contains structural variations,
then a mapper proceeds to split-mapping,
where it extracts the fixed length prefix and suffix from the read and query a separate short-length NavigaMer index.
In this manner, NavigaMer could support split mappings where a breakpoint is in the middle of a read and the read still contains reference-like prefix and suffix of decent lengths.
In the rare, unfortunate cases where NavigaMer completely fails,
it can always revert to a conventional read mapper for rescue mapping.
We intend to explore split mapping in our future work of expanding NavigaMer into a full mapper.

The source code of NavigaMer is available at: \url{https://anonymous.4open.science/r/NavigaMer-BF21/}.



% %%
% %% Print the bibliography
% %%
\bibliographystyle{plainnat}
\bibliography{oup-authoring-template/reference}

\end{document}
